% fig_throughput.tex — Зависимость времени обработки от объёма данных
\documentclass[tikz, border=2mm]{standalone}
\usepackage{fontspec}
\usepackage[russian]{babel}
\usepackage{pgfplots}
\pgfplotsset{compat=1.18}
\setsansfont{PT Sans}
\pgfplotsset{
    every axis/.append style={
        font=\small\sffamily,
        width=12cm,
        height=7cm
    },
    every tick label/.append style={font=\scriptsize\sffamily}
}

\begin{document}
\begin{tikzpicture}
\begin{axis}[
    xlabel={Объём данных, тыс. событий},
    ylabel={Время обработки, с},
    xmin=0, xmax=100,
    ymin=0, ymax=25,
    xtick={0,20,40,60,80,100},
    ytick={0,5,10,15,20,25},
    grid=both,
    grid style={line width=.1pt, draw=gray!30},
    major grid style={line width=.2pt, draw=gray!50},
    minor tick num=1,
    legend style={
        at={(0.5,-0.25)},
        anchor=north,
        legend columns=2,
        font=\small\sffamily
    },
    legend cell align={left},
    tick label style={/pgf/number format/use comma},
    table/col sep=semicolon
]

% Данные из inside-zi-2026.pdf, с. 4: время обработки атрибутивного метаграфа
\addplot[
    color=blue!80!black,
    mark=*,
    thick,
    mark size=2.0pt
]
coordinates {
    (10, 1.2)
    (20, 2.5)
    (30, 4.1)
    (40, 6.0)
    (50, 8.3)
    (60, 11.0)
    (70, 14.2)
    (80, 17.8)
    (90, 21.5)
    (100, 24.8)
};

\addlegendentry{AMGD (Разработанная методика)}

% Для сравнения: традиционный SIEM-подход (оценка по данным 10_Латыпов ИТ.pdf)
\addplot[
    color=red!80!black,
    mark=square*,
    thick,
    mark size=2.0pt
]
coordinates {
    (10, 2.0)
    (20, 4.8)
    (30, 8.5)
    (40, 13.0)
    (50, 18.0)
    (60, 24.0)
    % выше 60 тыс. — резкий рост, но обрежем для масштаба
};

\addlegendentry{SIEM-правила}

\end{axis}
\end{tikzpicture}
\end{document}